\documentclass[../../../include/open-logic-section]{subfiles}

\begin{document}

\olfileid[hi]{his}{set}{pathology}
\olsection{विकृतियाँ}

किंतु \emph{सीमा} की परिभाषा से कुछ काफ़ी ``विकृत'' रचनाएँ संभव निकलीं।

लगभग 1830 के दशक में बोल्ज़ानो ने ऐसा फलन खोजा जो \emph{सर्वत्र सतत}, किंतु
\emph{कहीं भी अवकलनीय नहीं} था। (दुर्भाग्य से बोल्ज़ानो ने इसे कभी प्रकाशित
नहीं किया; गणितज्ञों का इस विचार से पहला सामना 1872 में हुआ, जब
वाइयरश्ट्रास ने स्वतंत्र रूप से यही विचार खोजा।)\footnote{इस इतिहास का अत्यंत
विस्तृत प्रलेखन \href{http://en.wikipedia.org/wiki/Weierstrass_function}{वाइयरश्ट्रास
फलन} पर विकिपीडिया लेख की पाद-टिप्पणियों में है।} कम-से-कम कहें तो यह काफ़ी
आश्चर्यजनक था। $|x|$ जैसे फलन खोजना सरल है जो सर्वत्र सतत हों किंतु किसी
विशेष बिंदु पर अवकलनीय न हों। किंतु कोई फलन जो सर्वत्र सतत हो पर \emph{कहीं
भी} अवकलनीय न हो, बिल्कुल भिन्न जीव है। एक क्षण सोचें कि आप ऐसा फलन कैसे
खींचने का प्रयास करेंगे। उसे सतत सुनिश्चित करने के लिए आपको पृष्ठ से कलम
कभी हटाए बिना उसे खींच पाना चाहिए; किंतु उसे कहीं भी अवकलनीय न होना सुनिश्चित
करने के लिए आपको निरंतर अपनी कलम की दिशा अचानक बदलनी होगी।

इसके बाद और ``विकृतियाँ'' आईं। 5 जनवरी 1874 को कांटोर ने डेडेकिंड को पत्र
लिखकर यह समस्या रखी:
\begin{quote}
क्या किसी पृष्ठ (मान लें उसकी सीमा सहित कोई वर्ग) को किसी रेखा (मान लें उसके
अंतबिंदुओं सहित कोई सीधी रेखा) से एकैकी रूप से इस तरह सहसंबद्ध किया जा सकता
है कि पृष्ठ के प्रत्येक बिंदु के अनुरूप रेखा का एक बिंदु हो और इसके विपरीत
रेखा के प्रत्येक बिंदु के अनुरूप पृष्ठ का एक बिंदु हो?

इस समय भी मुझे लगता है कि इस प्रश्न का उत्तर बहुत कठिन है---यद्यपि यहाँ भी
\emph{नहीं} कहने की प्रेरणा इतनी प्रबल होती है कि कोई प्रमाण को लगभग अनावश्यक
मानना चाहेगा। [\citealt{Gouvea2011} में उद्धृत]
\end{quote}
किंतु 1877 में कांटोर ने सिद्ध किया कि वे गलत थे। वास्तव में किसी रेखा और
किसी वर्ग पर बिंदुओं की संख्या ठीक समान है। 29 जून 1877 को उन्होंने डेडेकिंड
को लिखा, ``\emph{je le vois, mais je ne le crois pas}''; अर्थात् ``मैं इसे
देखता हूँ, पर विश्वास नहीं करता''। गणित के ``स्वीकृत इतिहास'' में इसे प्रायः
इस संकेत के रूप में लिया जाता है कि ये नए परिणाम तत्कालीन गणितज्ञों को कितने
\emph{शब्दशः अविश्वसनीय} लगे थे। (पत्राचार \citet{Gouvea2011} में प्रस्तुत
है और हम \olref[his][set][mythology]{sec} में इस पर लौटेंगे। कांटोर के प्रमाण
की रूपरेखा \olref[his][set][cantorplane]{sec} में है।)

कांटोर के परिणाम से प्रेरित होकर पियानो ने विचार करना आरंभ किया कि क्या किसी
रेखा को किसी समतल पर \emph{सुचारु रूप से} प्रतिचित्रित करना संभव है। यह एक
\emph{अंतरिक्ष-पूर्ति वक्र} होता। \citeyear{Peano1890} में पियानो ने ठीक ऐसा
वक्र बनाया। यह सच में अंतर्ज्ञान-विरोधी है: यूक्लिड ने रेखा को ``चौड़ाई-रहित
लंबाई'' (पुस्तक I, परिभाषा 2) कहा था, किंतु पियानो ने दिखाया कि रेखा को
उपयुक्त रूप से मोड़कर उसकी लंबाई को चौड़ाई में बदला जा सकता है।
\citeyear{Hilbert1891} में हिल्बर्ट ने कुछ अधिक अंतर्ज्ञान-सुलभ
अंतरिक्ष-पूर्ति वक्र का वर्णन और उसे स्पष्ट करने वाले कुछ चित्र दिए। वक्र का
निर्माण क्रमशः होता है और निर्माण के पहले छह चरण ये हैं:
\begin{center}
\begin{tikzpicture}
\begin{scope}[xshift=20pt, yshift=20pt]
	\draw [oldiagcolorC, l-system={Hilbert curve, step=40pt, angle=90, axiom=L, order=1}] lindenmayer system;
\end{scope}
\begin{scope}[xshift=110pt, yshift=10pt]
	\draw [oldiagcolorC, l-system={Hilbert curve, step=20pt, angle=90, axiom=L, order=2}] lindenmayer system;
\end{scope}
\begin{scope}[xshift=205pt, yshift=5pt]
	\draw [oldiagcolorC, l-system={Hilbert curve, step=10pt, angle=90, axiom=L, order=3}] lindenmayer system;
\end{scope}
\begin{scope}[xshift=2.5pt, yshift=-97.5pt]
	\draw [oldiagcolorC, l-system={Hilbert curve, step=5pt, angle=90, axiom=L, order=4}] lindenmayer system;
\end{scope}
\begin{scope}[xshift=101.25pt, yshift=-98.75pt]
	\draw [oldiagcolorC, l-system={Hilbert curve, step=2.5pt, angle=90, axiom=L, order=5}] lindenmayer system;
\end{scope}
\begin{scope}[xshift=200.625pt, yshift=-99.375pt]
	\draw [oldiagcolorC, l-system={Hilbert curve, step=1.25pt, angle=90, axiom=L, order=6}] lindenmayer system;
\end{scope}
\draw[gray] (0pt,0pt) rectangle (80pt, 80pt);
\draw[gray] (100pt,0pt) rectangle (180pt, 80pt);
\draw[gray] (200pt,0pt) rectangle (280pt, 80pt);
\draw[gray] (0pt,-100pt) rectangle (80pt, -20pt);
\draw[gray] (100pt,-100pt) rectangle (180pt, -20pt);
\draw[gray] (200pt,-100pt) rectangle (280pt, -20pt);
\end{tikzpicture}  
\end{center}
सीमा में---एक संकल्पना जिसे अब तक कठोर परिभाषा मिल चुकी थी---पूरा वर्ग ठोस
!!{colorC} रंग से भर जाता है। और प्रसंगवश, हिल्बर्ट का वक्र सर्वत्र सतत है
पर कहीं भी अवकलनीय नहीं; अंतर्ज्ञानतः इसलिए कि अनंत सीमा में फलन हर क्षण
अचानक दिशा बदलता है। (हम \olref[his][set][hilbertcurve]{sec} में हिल्बर्ट
के निर्माण की अधिक विस्तृत रूपरेखा देंगे।)

अच्छा हो या बुरा, इन ``विकृत'' ज्यामितीय रचनाओं को ज्यामितीय अंतर्ज्ञान के
आह्वान पर संदेह करने का कारण माना गया। वे गणित करने के उस नए ढंग के लिए
लगभग \emph{प्रचार} बन गईं जिसकी परिणति समुच्चय सिद्धांत में होनी थी। विषय के
बाद के मिथक-निर्माण में बार-बार दोहराया गया कि ये परिणाम पूर्णतः कठोर और
पूर्णतः चौंकाने वाले, दोनों थे। इस तरह उन्होंने दोहरा उद्देश्य पूरा किया:
ज्यामितीय अंतर्ज्ञान पर भरोसा करने के विरुद्ध चेतावनी और सोचने के नए ढंगों
की उर्वरता का प्रदर्शन।

\end{document}
